\section{Exceptional terms and the power-saving induction}
\label{sec:induction}

We now combine the scale comparison, which applies to every small-residue
assignment, with the cancellation obtained for good assignments. For good
assignments the nonexceptional contribution is already negligible; we control
the exceptional parts by expanding the prime measures. The other assignments
are bounded directly by the scale comparison. Expanding measures rather than
comparing integrands pointwise keeps the argument valid for signed Fourier
lifts throughout.

The finite set $V$ and the constants $P,M_0,\beta,\gamma,\tau,\sigma_0$
have now been fixed. Lemma~\ref{lem:good-residues} has also established the
fixed positive proportion $\rho$ of good small-residue assignments.
The prime threshold $P$ was chosen so that
\begin{equation}
 \prod_{p\ge P}(1+p^{-3})-1\le\frac14.
 \label{eq:exceptional-euler-budget}
\end{equation}
For every positive integer $n$, define
\begin{equation}
 \mathcal Y(n)=
 \max_{\substack{A\subseteq[n]\\ A\text{ square-difference-free}}}Y(n,A).
 \label{eq:maximal-functional}
\end{equation}
This is a finite nonnegative number, by Proposition~\ref{prop:signed-holder}.
Write $\eta_p=\varepsilon_p(\F_p^V)$ for the mass of the exceptional measure,
so $0\le \eta_p\le p^{-4}$.  For
$B\subseteq\mathcal P_N$, put
\begin{equation}
 q_B=\prod_{p\in B}p,\qquad
 \eta_B=\prod_{p\in B}\eta_p,\qquad
 n_B=\left\lceil\frac{N}{(M_0q_B)^2}\right\rceil.
 \label{eq:exceptional-subset-notation}
\end{equation}
Empty products have value $1$, so $q_\varnothing=\eta_\varnothing=1$ and
$n_\varnothing=\lceil N/M_0^2\rceil$.

\begin{proposition}\label{prop:induction-recurrence}
There are absolute constants $K<\infty$ and $N_1\ge2$ such that, for every
integer $N\ge N_1$,
\begin{align}
 \mathcal Y(N)
 &\le K N^{-\sigma_0}
 +(1+u_N)\left((1-\rho)\mathcal Y(n_\varnothing)
       +\rho\!
       \sum_{\substack{\varnothing\ne B\subseteq\mathcal P_N\\q_B\le N^\tau}}
            \eta_B\mathcal Y(n_B)\right),
 \label{eq:functional-recurrence}\\
 u_N&=(1+M_0^2N^{2\tau-1})^d-1.
 \label{eq:uniform-ceiling-error}
\end{align}
In particular, $u_N\to0$ as $N\to\infty$.  The constants $K,N_1$ do not
depend on any proposed induction bound for $\mathcal Y$.
\end{proposition}

\begin{proof}
Fix a square-difference-free set $A\subseteq[N]$.  Write
\[
 G_s=\mathcal L_{[N],\mathcal P_N,Q}f_s,
 \qquad
 \Phi_{\mathbf s}((z_v)_{v\in V})=\prod_{v\in V}G_{s_v}(z_v),
\]
where $\mathbf s=(s_v)_{v\in V}$ is a small-residue assignment.  By
\eqref{eq:small-residue-input}, $1_A=\E_s f_s$.  The lift is linear, and the
product-space integral is multilinear, so
\begin{equation}
 Y(N,A)=\E_{\mathbf s}
       \int\Phi_{\mathbf s}\,d\bigotimes_{p\in\mathcal P_N}\mu_p.
 \label{eq:residue-average-functional}
\end{equation}

For $B\subseteq\mathcal P_N$, define the nonnegative finite measure
\[
 \kappa_B=
   \bigotimes_{p\in B}\varepsilon_p
   \otimes\bigotimes_{p\in\mathcal P_N\setminus B}\mu_p.
\]
Its mass is $\eta_B$.  The finite product expansion gives the following
identity of signed measures:
\begin{equation}
 \bigotimes_{p\in\mathcal P_N}\mu_p
 -\bigotimes_{p\in\mathcal P_N}(\mu_p-\varepsilon_p)
 =\sum_{\varnothing\ne B\subseteq\mathcal P_N}
       (-1)^{|B|+1}\kappa_B.
 \label{eq:signed-measure-inclusion-exclusion}
\end{equation}
For any real integrand $\Phi$, integration followed by the triangle inequality
therefore yields
\begin{equation}
 \int\Phi\,d\bigotimes_p\mu_p
 \le
 \left|\int\Phi\,d\bigotimes_p(\mu_p-\varepsilon_p)\right|
 +\sum_{B\ne\varnothing}\left|\int\Phi\,d\kappa_B\right|.
 \label{eq:exceptional-signed-bound}
\end{equation}
This step makes no sign assumption on $\Phi$.  Notice also that the measure
outside $B$ in each summand is the full measure $\mu_p$.

We first bound the summands with $q_B>N^\tau$.  A summand of zero mass is
zero, so assume $\eta_B>0$.  Each normalized exceptional measure
$\varepsilon_p/\eta_p$ has uniform single-label marginals: its invariance under
simultaneous translation implies that any one label is uniformly distributed
on $\F_p$.  The same uniformity holds for $\mu_p$.  Consequently every
single-slot marginal of $\kappa_B/\eta_B$ is uniform on
$X_{\mathcal P_N}$.  No independence between different slots is needed.
Ordinary H\"older's inequality gives
\begin{equation}
 \left|\int\Phi_{\mathbf s}\,d\kappa_B\right|
 \le\eta_B\prod_{v\in V}\|G_{s_v}\|_{L^d(X_{\mathcal P_N})}
 \le C_\epsilon\eta_BN^\epsilon
 \qquad(\epsilon>0).
 \label{eq:exceptional-holder-moment}
\end{equation}
The last inequality is Lemma~\ref{lem:lift-moments}, applied at the
fixed even order $d$ and with the loss in each factor chosen sufficiently
small.  The bound $|f_s|\le M_0$ only introduces a fixed factor.  The constant
$C_\epsilon$ is independent of $B$, $A$, $N$, and $\mathbf s$.

The exceptional masses are summable with the extra factor $q_B$:
\begin{equation}
 \sum_{\varnothing\ne B\subseteq\mathcal P_N}q_B\eta_B
 =\prod_{p\in\mathcal P_N}(1+p\eta_p)-1
 \le\prod_{p\ge P}(1+p^{-3})-1\le\frac14.
 \label{eq:weighted-exceptional-mass}
\end{equation}
It follows from \eqref{eq:exceptional-holder-moment} that
\[
 \sum_{\substack{\varnothing\ne B\subseteq\mathcal P_N\\q_B>N^\tau}}
   \left|\int\Phi_{\mathbf s}\,d\kappa_B\right|
 \le\frac{C_\epsilon}{4}N^{-\tau+\epsilon}.
\]
Choose $\epsilon=\tau/2$.  Since $\sigma_0=\gamma\tau/4<\tau/2$, we obtain
\begin{equation}
 \sum_{q_B>N^\tau}
   \left|\int\Phi_{\mathbf s}\,d\kappa_B\right|
 \le K_1N^{-\sigma_0},
 \label{eq:large-exceptional-tail}
\end{equation}
uniformly in the residue assignment.  Thus the number of exceptional subsets
has already been accounted for by the convergent product
\eqref{eq:weighted-exceptional-mass}.

Now suppose $q_B\le N^\tau$ and $\eta_B>0$.  Condition on all the labels at
primes in $B$ under $\kappa_B/\eta_B$.  The remaining prime law is exactly
$\bigotimes_{p\in\mathcal P_N\setminus B}\mu_p$.  For each realization, the
function in slot $v$ is $G_{s_v}$ specialized at its prescribed labels in $B$.
The mixed assertion of Proposition~\ref{prop:scale-transfer} applies to these
arbitrary slotwise residues and labels.  With
\[
 D_B=M_0q_B,\qquad \lambda_B=\frac{D_B^2n_B}{N},
\]
it gives
\begin{equation}
 \left|\int\Phi_{\mathbf s}\,d\kappa_B\right|
 \le\eta_B\left(\lambda_B^d\mathcal Y(n_B)+K_2N^{-\sigma_0}\right).
 \label{eq:small-exceptional-transfer}
\end{equation}
The same inequality is trivially true when $\eta_B=0$.  Crucially, the error
and its lower threshold for $N$ are uniform in $B$ and in every conditioned
tuple.  No induction hypothesis has been used.

The ceiling in the smaller scale produces a uniform multiplicative error.
For $q_B\le N^\tau$,
\[
 1\le\lambda_B\le1+\frac{D_B^2}{N}
       \le1+M_0^2N^{2\tau-1},
\]
and hence $\lambda_B^d\le1+u_N$.  Moreover,
$\sum_{B\ne\varnothing}\eta_B\le1/4$ by
\eqref{eq:weighted-exceptional-mass}.  Summing the errors in
\eqref{eq:small-exceptional-transfer} is therefore harmless.

If $\mathbf s$ is good, Proposition~\ref{prop:nonexceptional-contraction}
bounds the first term on the right of \eqref{eq:exceptional-signed-bound} by
$K_3N^{-\sigma_0}$.  Combining this with
\eqref{eq:large-exceptional-tail} and
\eqref{eq:small-exceptional-transfer}, we obtain
\[
 \int\Phi_{\mathbf s}\,d\bigotimes_p\mu_p
 \le K_4N^{-\sigma_0}
 +(1+u_N)
   \sum_{\substack{\varnothing\ne B\subseteq\mathcal P_N\\q_B\le N^\tau}}
      \eta_B\mathcal Y(n_B).
\]
For a nongood $\mathbf s$, apply the mixed scale comparison directly with
$B=\varnothing$:
\[
 \left|\int\Phi_{\mathbf s}\,d\bigotimes_p\mu_p\right|
 \le(1+u_N)\mathcal Y(n_\varnothing)+K_2N^{-\sigma_0}.
\]
Averaging these two estimates in
\eqref{eq:residue-average-functional} gives
\eqref{eq:functional-recurrence}, since the good assignments have probability
$\rho$.  The estimate is uniform in $A$, so we may take its maximum.
Finally, $2\tau<1$ by the fixed parameter choice, and $M_0,d$ are fixed, so
$u_N\to0$.  All constants in the proof came from the preceding estimates and
fixed parameters; none depended on an induction constant.
\end{proof}

\begin{theorem}\label{thm:functional-decay}
There are absolute constants $a_0>0$ and $C_0<\infty$ such that, for every
positive integer $N$ and every square-difference-free set $A\subseteq[N]$,
\[
 Y(N,A)\le C_0N^{-a_0}.
\]
\end{theorem}

\begin{proof}
We make all remaining parameter choices before starting the induction.
The finite construction fixes $h,t_0,\ell,d$; the conditional and kernel estimates
fix $\gamma$ and $\beta$.  The threshold $P$ was then chosen, and with it
$M_0$, $\tau$, and $\sigma_0$.  In particular the probability $\rho$ of
Lemma~\ref{lem:good-residues} is already a fixed positive number.  We shall
not enlarge $P$ after choosing the exponent below.

Put $r=1-3\rho/4$, so $1/4\le r<1$, and choose
\begin{equation}
 a_0=\min\left\{\frac14,\frac{\sigma_0}{2},
                  \frac{\log(1/r)}{4\log M_0}\right\}.
 \label{eq:absolute-exponent-choice}
\end{equation}
Table~\ref{tab:parameter-roles} summarizes the roles of the parameters;
among the listed quantities, only those in the last row vary with
the interval length.
\begin{table}[htbp]
\centering
\small
\begin{tabular}{@{}p{.20\textwidth}p{.71\textwidth}@{}}
\toprule
Parameters & Role and order of choice\\
\midrule
$h,t_0,\ell,d$ & Fixed cycle length, marking depth, block count, and tuple size
(Section~\ref{sec:tuple-laws}).\\[3pt]
$\gamma,\beta$ & Fixed conductor-decay and rational-frequency cutoff exponents
(Sections~\ref{sec:functional} and~\ref{sec:square-kernel}).\\[3pt]
$P,M_0$ & Prime-law threshold and modulus containing all smaller primes;
fixed before choosing the decay exponent
(Section~\ref{sec:scale-transfer}).\\[3pt]
$\tau,\sigma_0$ & Fixed partition exponent and common error exponent;
chosen in~\eqref{eq:scale-parameters}.\\[3pt]
$\rho,a_0$ & Fixed positive residue proportion and induction exponent;
$a_0$ is chosen only after $M_0$ and $\rho$ in~\eqref{eq:absolute-exponent-choice}.\\[3pt]
$Q,H,\mathcal P_N$ & Scale-dependent cutoffs and prime set, defined from $N$
in~\eqref{eq:scale-parameters}.\\
\bottomrule
\end{tabular}
\caption{Parameter roles. Every construction size and exponent is fixed
independently of $N$ and $A$; the Fourier cutoffs and available prime
coordinates change with $N$.}
\label{tab:parameter-roles}
\end{table}


Every quantity in~\eqref{eq:absolute-exponent-choice} is fixed and positive. Thus
$0<a_0<\min(1/2,\sigma_0)$, and
\begin{equation}
 \theta:=M_0^{2a_0}(1-\rho+\rho/4)
          =M_0^{2a_0}r\le\sqrt r<1.
 \label{eq:strict-induction-contraction}
\end{equation}
There is no requirement that $a_0$ be numerically useful.

For every finite $\mathcal P_N$, the bound $\eta_p\le p^{-4}$ and
$2a_0\le1$ give
\begin{equation}
 \sum_{\varnothing\ne B\subseteq\mathcal P_N}\eta_Bq_B^{2a_0}
 =\prod_{p\in\mathcal P_N}(1+\eta_p p^{2a_0})-1
 \le\prod_{p\ge P}(1+p^{-3})-1\le\frac14.
 \label{eq:induction-weighted-mass}
\end{equation}

Let $\kappa=(1-\theta)/2>0$.  Choose an integer $N_*\ge N_1$ so large that,
for every $N\ge N_*$,
\begin{equation}
 \theta u_N\le\kappa,\qquad
 K N^{a_0-\sigma_0}\le\kappa,
 \label{eq:induction-threshold-choice}
\end{equation}
where $K,N_1$ are supplied by
Proposition~\ref{prop:induction-recurrence}.  This is possible because
$u_N\to0$ and $a_0<\sigma_0$.  The choice of $N_*$ is independent of $C_0$.
Choose $C_0\ge1$ large enough that
\begin{equation}
 \mathcal Y(n)\le C_0 n^{-a_0}\qquad(1\le n<N_*).
 \label{eq:finite-induction-base}
\end{equation}
Only finitely many values are involved, so such a finite constant exists.
For an elementary bound, the lift in $Y(n,A)$ has at most $2n^{2\beta}$
frequencies and each coefficient has modulus at most $1$.  Since its tuple
law is a probability measure,
$0\le Y(n,A)\le(2n^{2\beta})^d$.  Thus, for example,
$C_0=\max\{1,2^dN_*^{2\beta d+a_0}\}$ suffices in
\eqref{eq:finite-induction-base}.  This also covers $n=1$ and the cases of
an empty prime set.

Suppose now that $N\ge N_*$ and
$\mathcal Y(n)\le C_0n^{-a_0}$ for every positive integer $n<N$.
Every smaller scale appearing in \eqref{eq:functional-recurrence} lies in
this range.  Indeed, $M_0q_B\ge8$ gives
\[
 1\le n_B\le N/64+1<N\qquad(N\ge2).
\]
The ceiling goes in the favorable direction for a negative exponent:
\[
 n_B^{-a_0}
 \le N^{-a_0}M_0^{2a_0}q_B^{2a_0}.
\]
Apply the induction hypothesis at every such scale in
\eqref{eq:functional-recurrence}, and then use
\eqref{eq:induction-weighted-mass}.  Extending the positive subset sum to all
nonempty $B$ can only increase the bound, and therefore
\begin{align*}
 \mathcal Y(N)
 &\le K N^{-\sigma_0}
 +(1+u_N)C_0N^{-a_0}M_0^{2a_0}
       \left(1-\rho+\rho
           \sum_{B\ne\varnothing}\eta_Bq_B^{2a_0}\right)\\
 &\le K N^{-\sigma_0}+(1+u_N)\theta C_0N^{-a_0}\\
 &\le \kappa N^{-a_0}+(1-\kappa)C_0N^{-a_0}
 \le C_0N^{-a_0}.
\end{align*}
In the third line we used
\eqref{eq:induction-threshold-choice} and
$\theta+\kappa=1-\kappa$; in the last line we used $C_0\ge1$.
This closes strong induction.  In particular, increasing $C_0$ to cover the
finite base range does not change the validity of the large-scale step.
\end{proof}

\begin{proof}[Proof of Theorem~\ref{thm:main}]
Let $N\ge1$ and let $A\subseteq[N]$ have no positive integer-square
difference.  Put
$G=\mathcal L_{[N],\mathcal P_N,Q}1_A$.  Its constant Fourier coefficient is
$|A|/N$, and every nonconstant character on $X_{\mathcal P_N}$ has uniform
mean zero.  The mean inequality in Lemma~\ref{lem:functional-properties} and
Theorem~\ref{thm:functional-decay} yield
\[
 \frac{|A|}{N}=|\E G|
 \le Z_{\mathcal P_N}(G)^{1/d}
 =Y(N,A)^{1/d}
 \le C_0^{1/d}N^{-a_0/d}.
\]
Thus $C=C_0^{1/d}<\infty$ and $c=a_0/d>0$ are absolute constants for which
$|A|\le CN^{1-c}$ for every such $N$ and $A$.
\end{proof}
